\documentclass[12pt]{article}
\usepackage[utf8]{inputenc}
\usepackage[T1]{fontenc}
\usepackage[margin=1in]{geometry}
\usepackage{setspace}
\usepackage{xcolor}
\usepackage[colorlinks=true, urlcolor=blue!60!black, linkcolor=blue!60!black]{hyperref}
\setlength{\parindent}{0em}
\setlength{\parskip}{0.6em}
\renewcommand{\baselinestretch}{1.0}
\begin{document}
\begin{center}
{\LARGE \textit{Stuart v. Laird}, 5 U.S. (1 Cranch) 299 (1803)\par}
\vspace{0.5em}
{\large Case Brief --- Two-Page Edition\par}
\end{center}

\section{Memory Jogger}
Days after \textit{Marbury}, the Court accepted the Republicans' repeal of the 1801 Judiciary Act: Congress may transfer cases between lower courts, and practice settled circuit riding (\href{https://supreme.justia.com/cases/federal/us/5/299/}{Justia}).

\section{Detailed Case Facts}
John Laird's debt judgment against Hugh Stuart, won in a circuit court created by the Federalists' 1801 ``Midnight Judges Act,'' was transferred to a reconstituted circuit court after the Republican Congress repealed that Act in 1802 and restored circuit riding by Supreme Court Justices. Chief Justice Marshall, riding circuit, entered judgment on Stuart's forthcoming bond.

\section{Procedural History}
Stuart brought a writ of error challenging the transfer and the Justices' authority to sit on circuit without separate commissions.

\section{Judicial Votes}
Affirmed 5--0. Paterson wrote for Cushing, Chase, Washington, and Moore; Marshall did not participate.

\section{Holding}
Congress may create inferior courts and transfer pending cases among them. Long practice and acquiescence settled that Justices may ride circuit under their existing commissions. The Court did not address whether abolishing the sixteen judgeships violated the Good Behavior Clause.

\section{Analysis of Opinions}
\textbf{Paterson:} the power to establish inferior tribunals includes the power to transfer causes, and the bond was ``an appendage to the cause.'' Circuit riding since 1789 was a ``contemporary interpretation of the most forcible nature.'' \textbf{No separate opinions:} despite private doubts by Marshall and Chase, the Justices closed ranks amid threats of impeachment.

\section{Examples: Future Applications}
\textbf{Same-Side:} splitting the Ninth Circuit and transferring pending appeals. \textbf{Opposite-Side:} abolishing one judge's seat to oust her violates the Good Behavior Clause. \textbf{Fence-Sitter:} eighteen-year active-service limits for Justices.

\section{Critique}
Liquidation theorists see \textit{Stuart} as proof that early practice can settle text; critics note that thirteen years of protested practice overrode textual objections and that the tenure question was dodged. Others praise it as statecraft paired with \textit{Marbury}.

\section{Key Quotations}
\begin{quote}
``To this objection, which is of recent date, it is sufficient to observe, that practice and acquiescence under it for a period of several years, commencing with the organization of the judicial system, affords an irresistible answer, and has indeed fixed the construction.''
\end{quote}
\noindent --- Justice Paterson, for the Court, 5 U.S. (1 Cranch) at 309.

\end{document}
